\documentclass{article}

\usepackage{graphicx,amsmath,amssymb}

\title{量子力学}
\author{a}
\date{April 2026}

\begin{document}

\maketitle

\section{歴史系背景と問題意識 - 古典物理学の限界}

19世紀末までに、ニュートン力学とマクスウェル電磁気学によって、物理現象の大部分は説明されたと考えられていた。しかし、この理論には明確な限界があった。特に、微視的現象においては、古典理論の予測と実験結果との間に重大な不一致が現れた。
\begin{itemize}
\item{黒体放射と紫外線発散}: 
黒体放射のスペクトルは古典論では説明できず、高周波領域でエネルギーが発散するという「紫外線発散」が現れる。プランクは、エネルギーが連続ではなく離散的に交換されると仮定することでこの問題を解決した。

\item{光電効果}: 
光電効果では、光の強度ではなく振動数が電子放出の有無を決める。アインシュタインは、光が粒子的性質を持つこと、すなわちエネルギー量子として振る舞うことを提案した。

\item{原子の安定性の問題}: 
古典電磁気学では、電子は加速運動により電磁波を放出し、最終的には原子核に落ち込むはずである。しかし実際には原子は安定であり、この事実を古典理論では説明できない。
\end{itemize}
以上の事情から、古典力学における状態の記述および時間発展の枠組みは、そのままでは微視的世界に適用できないことが明らかとなる。状態、観測量、および時間発展の法則を再構築する新理論の必要が明らかとなった。

\section{記法と量子力学の公理}
量子力学にはいくつかの原理がある。歴史的にはこれは実験結果を説明するために段階的に見出されてきたものであるが、話の流れのために先に第一原理を導入してしまうこととする。
本講義ノートでは、量子力学の理論体系を以下の公理的枠組みに基づき記述する。

\subsection{状態（公理I）}
系の物理的状態は、複素ヒルベルト空間の元（状態ベクトル）$\psi$ によって表される。$\psi$が複素係数倍（全体位相）の違いを除いて同一なら、物理的に同一の状態である。ノルムが1になるような係数を選ぶことを「正規化」という。

\subsection{観測量（公理II）}
物理量はヒルベルト空間上の自己共役演算子 $\hat{A}$ によって表される。本ノートでは、演算子であることを明示するためにハット（$\hat{\quad}$）を付す。

\subsection{測定と確率（公理III）}
状態 $\psi$（ただし正規化されているとする）において物理量 $\hat{A}$ を測定したとき、得られる値は $\hat{A}$ の固有値のいずれかであり、その確率は状態ベクトルと固有状態の内積の絶対値の2乗（ボルンの規則）によって与えられる。

\subsection{時間発展（公理IV）}
時間発展は線形であり、かつ確率（内積）を保存するユニタリ変換として与えられる。

\section{単一粒子の非相対論的な量子力学}

\subsection{シュレーディンガー方程式の導出}
上の原理に従って、単一粒子の非相対論的な量子力学の基本方程式を考えてみる。
出発点は、古典力学における粒子のエネルギーと運動量の関係
\begin{equation}
    E = \frac{\mathbf{p}^2}{2m} + V(\mathbf{x})
\end{equation}
である。一方、ド・ブロイの関係式によれば、粒子は波としての性質を持ち、$E = \hbar \omega, \mathbf{p} = \hbar \mathbf{k}$ を満たす。
粒子は波として表現できるのであるから、系の状態を表すヒルベルト空間を $\mathcal{H}:=C^\infty[\mathbb{R}^3, \mathbb{C}]$ としてみよう。
その基本的な例として平面波 $\psi(\mathbf{x}, t) = e^{i(\mathbf{k} \cdot \mathbf{x} - \omega t)}$ を考える。これに時間微分および空間微分を作用させると、以下のようにエネルギーや運動量を取り出すことができる。
\begin{equation}
    i\hbar \frac{\partial}{\partial t} \psi = \hbar \omega \psi = E \psi, \quad -i\hbar \nabla \psi = \hbar \mathbf{k} \psi = \mathbf{p} \psi
\end{equation}
が成り立つ。したがって、エネルギーと運動量にはそれぞれ
\begin{equation}
    \hat{E} = i\hbar \frac{\partial}{\partial t}, \quad \hat{\mathbf{p}} = -i\hbar \nabla
\end{equation}
という演算子を対応させるのが自然である。これが対応原理である。

そこで、先ほどの古典的関係にこれらの演算子を代入すると、シュレーディンガー方程式
\begin{equation}
    i\hbar \frac{\partial}{\partial t} \psi = \left( -\frac{\hbar^2}{2m} \nabla^2 + \hat{V}(\mathbf{x}) \right) \psi = \hat{H} \psi
\end{equation}
が得られる。すなわち、古典的なエネルギー保存則を波動関数に作用する演算子の関係へと読み替えることで、量子力学の基本方程式が導かれる。これは線形な微分方程式であり、重ね合わせの原理を自然に満たしている。

シュレーディンガー方程式の重要な役割は、ポテンシャルと境界条件を与えたときに、許される状態とエネルギー固有値を具体的に決める点にある。典型的な例を三つ挙げる。
\subsection{シュレーディンガー方程式の基本的な解}
\paragraph{無限深井戸}
一次元で
\[
V(x)=
\begin{cases}
0 & (0 < x < L) \\
\infty & (\text{otherwise})
\end{cases}
\]
とする。壁の外では波動関数が消えなければならないため、定常状態 $\psi(x,t)=\phi(x)e^{-iEt/\hbar}$ の空間部分は
\[
\phi_n(x)=\sqrt{\frac{2}{L}}\sin\frac{n\pi x}{L}, \quad E_n=\frac{\hbar^2\pi^2 n^2}{2mL^2}\qquad(n=1,2,\dots)
\]
となる。境界条件だけでエネルギーが離散化されることが分かる。

\paragraph{有限ポテンシャル障壁}
たとえば
\[
V(x)=
\begin{cases}
V_0 & (0 < x < L) \\
0 & (\text{otherwise})
\end{cases}
\]
という障壁を考える。$E<V_0$ であっても、障壁内部では波動関数が指数関数的に減衰しながら残るため、反対側への透過確率は零にはならない。これは古典力学には現れないトンネル効果であり、ポテンシャルが有限であることの量子論的帰結である。

\paragraph{調和振動子}
ポテンシャル
\[
V(x)=\frac{1}{2}m\omega^2 x^2
\]
をもつ調和振動子では、正規化可能な固有状態を要求することで、エネルギー固有値は
\[
E_n = \hbar \omega \left(n + \frac{1}{2}\right)
\]
となる。基底状態のエネルギーが $0$ ではなく $\hbar\omega/2$ だけ残る点は、量子揺らぎの典型例である。

\subsection{交換関係とその意義}
位置演算子 $\hat{x}$ と運動量演算子 $\hat{p}$ は、正準交換関係
\[
[\hat{x}, \hat{p}] = i\hbar
\]
を満たす。位置表示では $\hat{x}=x$, $\hat{p}=-i\hbar \frac{d}{dx}$ であるから、任意の波動関数 $\psi(x)$ に対して
\[
[\hat{x},\hat{p}]\psi(x)
= x\left(-i\hbar \frac{d\psi}{dx}\right) + i\hbar \frac{d}{dx}(x\psi(x))
= i\hbar \psi(x)
\]
となり、交換関係が確かめられる。

この関係の意義は、位置と運動量が同時に鋭く定まる共通固有状態を持たないという点にある。さらに、運動量が空間並進の生成子であることや、後に述べる不確定性原理およびハイゼンベルク方程式の構造は、この交換関係に支えられている。

\subsection{不確定性原理}
状態 $\psi$ における期待値を $\langle \hat{A} \rangle$ とし、
\[
\Delta \hat{x} := \hat{x} - \langle \hat{x} \rangle, \quad \Delta \hat{p} := \hat{p} - \langle \hat{p} \rangle
\]
とおく。標準偏差は
\[
(\Delta x)^2 = \langle (\Delta \hat{x})^2 \rangle, \quad (\Delta p)^2 = \langle (\Delta \hat{p})^2 \rangle
\]
で与えられる。ここで、ベクトル $\Delta \hat{x}\psi$ と $\Delta \hat{p}\psi$ にコーシー--シュワルツの不等式を適用すると、
\[
(\Delta x)^2(\Delta p)^2 \ge \left| \langle \Delta \hat{x}\Delta \hat{p} \rangle \right|^2
\]
を得る。さらに、複素数 $z$ に対して $|z| \ge |\operatorname{Im} z|$ であることと、
\[
\operatorname{Im}\langle \Delta \hat{x}\Delta \hat{p} \rangle
= \frac{1}{2i}\langle [\Delta \hat{x}, \Delta \hat{p}] \rangle
= \frac{1}{2i}\langle [\hat{x}, \hat{p}] \rangle
= \frac{\hbar}{2}
\]
を用いれば、
\[
\Delta x \, \Delta p \ge \frac{\hbar}{2}
\]
が従う。したがって、不確定性原理は測定器の不完全さではなく、交換関係そのものから出てくる量子論の構造的な制約である。

\section{相対論的理論への動機}
非相対論的量子力学は多くの現象を見事に説明するが、その基本方程式
\[
i\hbar \frac{\partial}{\partial t}\psi = \hat{H}\psi
\]
は時間に関して1階、空間に関しては一般に2階であり、時間と空間を非対称に扱っている。この形式はガリレイ変換の下では自然だが、ローレンツ対称性とは両立しない。

さらに、ここまでの記述では粒子数を固定したヒルベルト空間を暗黙に前提していた。しかし、相対論と量子論を同時に考えると、十分高いエネルギーでは粒子の生成・消滅が起こりうるため、この前提自体が揺らぐ。

以降では、まず「一粒子の相対論的波動方程式」を作ろうとすると何が起こるかを見た上で、その限界が場の量子論を要請することを確認する。

\section{相対論的波動方程式}

\subsection{クライン--ゴルドン方程式の導出}
特殊相対性理論において、静止質量 $m$ をもつ自由粒子のエネルギー $E$ と運動量 $\mathbf{p}$ の間には以下の関係が成り立つ：
\begin{equation}
    E^2 - c^2 \mathbf{p}^2 = m^2 c^4
\end{equation}
このエネルギー・運動量関係に、非相対論的量子力学における対応原理（演算子への置き換え）を形式的に適用すると、複素関数 $\phi(x)$ に対するクライン--ゴルドン（KG）方程式が得られる。
\begin{equation}
    \left( \frac{1}{c^2} \frac{\partial^2}{\partial t^2} - \nabla^2 + \frac{m^2 c^2}{\hbar^2} \right) \phi(\mathbf{x}, t) = 0
\end{equation}

\subsection{負エネルギー解が棄却できない理由}
この関係からエネルギー $E$ を求めると $E = \pm \sqrt{c^2 p^2 + m^2 c^4}$ となり、負のエネルギー解が現れる。量子論では以下の理由により棄却不可能である。

\paragraph{1. 状態の完全性}
任意の波動関数を展開するための基底を形成するには、負エネルギー解もすべて含める必要がある。

\paragraph{2. 遷移の可能性}
正エネルギー状態から負エネルギー状態へ遷移する確率が有限となる。

\paragraph{3. 因果律と局所性}
因果律を守り、粒子の局所性を維持するためには、負エネルギー解との重ね合わせが数学的に必須となる。

\subsection{再解釈：一粒子から多粒子へ}
この「負エネルギー解の棄却不可能性」こそが、粒子数が不変であるという「一粒子力学」の限界を露呈させた。これを解決するには、負エネルギー状態を「反粒子」として解釈する場の量子論への移行が必要となる。

\section{ディラック方程式}

\subsection{ディラックの着想：KG方程式の因数分解}
ディラックは、クライン--ゴルドン方程式が「時間について2階微分」であるために確率密度の負値が生じると考えた。そこで、シュレーディンガー方程式と同じく時間について1階の微分方程式
\begin{equation}
    i\hbar \frac{\partial}{\partial t} \psi = \hat{H}_D \psi
\end{equation}
を目指した。相対論的な対称性（時間と空間の対等性）を保つためには、ハミルトニアン $\hat{H}_D$ は空間微分（運動量演算子）についても1階でなければならない。
そこで、ハミルトニアンを以下の線形結合として仮定する：
\begin{equation}
    \hat{H}_D = c (\alpha_1 \hat{p}_x + \alpha_2 \hat{p}_y + \alpha_3 \hat{p}_z) + \beta m c^2
\end{equation}

\subsection{ディラック代数の導出}
このハミルトニアンが相対論的なエネルギー・運動量の関係 $E^2 = c^2 p^2 + m^2 c^4$ を満たすためには、$\hat{H}_D^2$ がクライン--ゴルドン方程式の演算子と一致しなければならない。
実際に $\hat{H}_D^2$ を計算すると：
\begin{equation}
    \hat{H}_D^2 = c^2 \sum_{i,j} \frac{\alpha_i \alpha_j + \alpha_j \alpha_i}{2} \hat{p}_i \hat{p}_j + mc^3 \sum_i (\alpha_i \beta + \beta \alpha_i) \hat{p}_i + \beta^2 m^2 c^4
\end{equation}
これが $c^2 \sum \hat{p}_i^2 + m^2 c^4$ と一致するための条件として、係数 $\alpha_i, \beta$ が以下の反交換関係（ディラック代数）を満たすことが要求される：
\begin{align}
    \{ \alpha_i, \alpha_j \} &= 2\delta_{ij} \\
    \{ \alpha_i, \beta \} &= 0 \\
    \alpha_i^2 = \beta^2 &= 1
\end{align}
これらの関係は通常の複素数では満たされず、$\alpha_i, \beta$ は行列でなければならない。

\subsection{スピノールと4成分解釈}
上述の条件を満たす最小の次元は $4 \times 4$ 行列である。したがって、状態空間は複素数値の波動関数の空間から、4成分をもつ関数の空間へと拡張される。すなわち、波動関数 $\psi$ は単なるスカラー関数ではなく、4つの成分を持つ列ベクトル「スピノール」となる。
\begin{equation}
    \psi(x) = \begin{pmatrix} \psi_1(x) \\ \psi_2(x) \\ \psi_3(x) \\ \psi_4(x) \end{pmatrix}
\end{equation}
ここでいうスピンとは、粒子の位置や運動量とは独立な内部自由度であり、空間回転に対して状態がどのように変換するかを表す量である。この4つの成分は、「正エネルギーの2つの内部状態」と「負エネルギーの2つの内部状態」に対応しており、この内部自由度が後にスピン $1/2$ として解釈される。

\subsection{ディラック方程式の意義}
以上を踏まえると、ディラック方程式の意義は次の二点にまとめられる。
\begin{enumerate}
    \item \textbf{スピンの自然な導入}：電子が持つ回転に関する二値の内部自由度が、理論の数学的構造から自動的に現れ、それがスピン $1/2$ として解釈される。
    \item \textbf{反粒子の予言}：負エネルギー解の問題に対し、ディラックは「ディラックの海」という概念を提唱し、反粒子の存在を予言した。
\end{enumerate}
しかし、ここで得られたのはあくまで相対論的な一粒子方程式であり、粒子の生成・消滅や真空の構造を理論の内部で扱うには、なお場の量子論への移行が必要である。

\section{$n$ 粒子系の量子力学}

一粒子相対論の困難は、「粒子数を固定したまま量子論を作る」という前提そのものを見直す必要があることを示している。そこで次に、通常の量子力学が多粒子系をどのように記述しているかを整理し、その枠組みの限界を明確にする。

\subsection{固定粒子数の状態空間}
一粒子系では、状態は位置 $\mathbf{x}$ を変数にもつ波動関数 $\psi(\mathbf{x},t)$ で表された。これを $n$ 個の粒子へ拡張すると、状態は各粒子の位置
\[
(\mathbf{x}_1,\mathbf{x}_2,\dots,\mathbf{x}_n)
\]
を変数にもつ波動関数
\[
\Psi(\mathbf{x}_1,\mathbf{x}_2,\dots,\mathbf{x}_n,t)
\]
で表される。したがって、$n$ 粒子系のヒルベルト空間は一粒子空間のテンソル積
\[
\mathcal{H}_n = \mathcal{H}_1^{\otimes n}
\]
として構成される。非相対論的な $n$ 粒子ハミルトニアンの典型例は
\[
\hat{H}_n
= \sum_{i=1}^n \left( -\frac{\hbar^2}{2m_i}\nabla_i^2 + V_{\mathrm{ext}}(\mathbf{x}_i) \right)
+ \sum_{1\le i<j\le n} U(\mathbf{x}_i-\mathbf{x}_j)
\]
であり、外場との相互作用と粒子間相互作用を同時に含む。時間発展は
\[
i\hbar \frac{\partial}{\partial t}\Psi = \hat{H}_n \Psi
\]
で与えられる。

\subsection{同種粒子と量子統計}
粒子が区別可能であれば、各 $\mathbf{x}_i$ は「第 $i$ 粒子」の座標として扱える。このとき、$n$ 粒子波動関数全体は複素数値関数のベクトル空間をなし、粒子の交換
\[
s_{ij} : (\mathbf{x}_1,\dots,\mathbf{x}_i,\dots,\mathbf{x}_j,\dots,\mathbf{x}_n)
\mapsto
(\mathbf{x}_1,\dots,\mathbf{x}_j,\dots,\mathbf{x}_i,\dots,\mathbf{x}_n)
\]
は配置空間上の写像として定義できる。したがって、波動関数への作用
\[
(\hat{P}_{ij}\Psi)(\mathbf{x}_1,\dots,\mathbf{x}_n)
:=
\Psi(s_{ij}(\mathbf{x}_1,\dots,\mathbf{x}_n))
\]
を考えることができる。これは固定した写像 $s_{ij}$ との合成であるから、
\[
\hat{P}_{ij}(a\Psi+b\Phi)
=
a\,\hat{P}_{ij}\Psi+b\,\hat{P}_{ij}\Phi
\]
が成り立ち、$\hat{P}_{ij}$ は線形演算子である。

しかし、電子同士や光子同士のような同種粒子では、粒子のラベルそのものに物理的意味はない。したがって、この交換は物理状態を変えず、波動関数には高々全体位相しか掛からない。すなわち
\[
\hat{P}_{ij}\Psi = \lambda_{ij}\Psi,
\qquad |\lambda_{ij}|=1
\]
である。さらに、交換を二回行えば元に戻るので $\hat{P}_{ij}^2=1$ であり、
\[
\lambda_{ij}^2=1
\]
が従う。したがって $\lambda_{ij}=\pm 1$ である。ゆえに、同種粒子の波動関数は交換に対して
\[
\Psi(\dots,\mathbf{x}_i,\dots,\mathbf{x}_j,\dots)
= \pm
\Psi(\dots,\mathbf{x}_j,\dots,\mathbf{x}_i,\dots)
\]
を満たす。符号が $+$ の場合がボース粒子、$-$ の場合がフェルミ粒子である。特にフェルミ粒子では
\[
\Psi(\dots,\mathbf{x},\dots,\mathbf{x},\dots)=0
\]
となるため、同じ一粒子状態を二つ以上占められない。これがパウリの排他原理である。

\subsection{固定された $n$ の理論の限界}
このように、非相対論的量子力学の範囲では、粒子数 $n$ を固定したまま多体系を記述できる。原子、分子、固体中の電子系など、多くの問題はこの枠組みで出発できる。しかし、この形式では最初に粒子数 $n$ を決めてしまっており、その後の時間発展で粒子が生成・消滅する過程は表せない。

そこで、粒子数が $0,1,2,\dots$ と変わりうる状況を扱うには、各 $n$ 粒子空間をまとめた
\[
\mathcal{F}
= \bigoplus_{n=0}^{\infty} \mathcal{H}_n
\]
という直和空間を考える必要がある。これがフォック空間である。場の量子論は、この空間の上で生成・消滅演算子を導入し、粒子数が変化する力学を記述する理論として理解できる。

\section{ハイゼンベルク描像から場の量子論へ}

ここまでで、一粒子波動関数を基本変数とする記述は、非相対論的近似や固定粒子数の問題には有効でも、相対論的かつ粒子数可変の状況では根本的でないことが見えてきた。そこで、状態ベクトルの時間依存を追うよりも、時空の各点に割り当てられた演算子の振る舞いを前面に出す記述へ移る。

\subsection{時間発展演算子}
これまでの記述（シュレディンガー描像）では、演算子は時間に依存せず、状態ベクトル $\psi(t)$ が時間発展を担っていた。しかし、物理的に観測可能なのは期待値 $\langle \hat{A} \rangle = (\psi(t), \hat{A} \psi(t))$ であり、この時間依存性を演算子側に担わせても同じ物理を記述できる。

時間発展が連続で確率を保存するなら、状態は一パラメータのユニタリ演算子族 $\hat{U}(t)$ によって
\[
\psi(t) = \hat{U}(t)\psi(0), \quad \hat{U}(0)=1, \quad \hat{U}(t+s)=\hat{U}(t)\hat{U}(s)
\]
と書ける。微小時間 $\delta t$ に対しては
\[
\hat{U}(\delta t) = 1 - \frac{i}{\hbar}\hat{H}\delta t + O(\delta t^2)
\]
と展開でき、ユニタリ性から生成子 $\hat{H}$ は自己共役でなければならない。したがって $\hat{U}(t)$ は
\[
\frac{d}{dt}\hat{U}(t) = -\frac{i}{\hbar}\hat{H}\hat{U}(t)
\]
を満たす。ハミルトニアン $\hat{H}$ が時間に依存しない場合、この微分方程式の解は
\[
\hat{U}(t) = \exp\left( -\frac{i\hat{H}t}{\hbar} \right)
\]
である。

\subsection{ハイゼンベルクの方程式}
ハイゼンベルク描像では、状態ベクトルを $\psi_H = \psi(0)$ に固定し、演算子を
\[
\hat{A}_H(t) = \hat{U}^\dagger(t)\hat{A}_S(t)\hat{U}(t)
\]
と定義する。これを時間微分すると、
\begin{align*}
\frac{d}{dt}\hat{A}_H(t)
&= \frac{d\hat{U}^\dagger}{dt}\hat{A}_S\hat{U}
+ \hat{U}^\dagger \frac{\partial \hat{A}_S}{\partial t}\hat{U}
+ \hat{U}^\dagger \hat{A}_S \frac{d\hat{U}}{dt} \\
&= \frac{i}{\hbar}\hat{U}^\dagger \hat{H}\hat{A}_S\hat{U}
+ \hat{U}^\dagger \frac{\partial \hat{A}_S}{\partial t}\hat{U}
- \frac{i}{\hbar}\hat{U}^\dagger \hat{A}_S\hat{H}\hat{U}
\end{align*}
となる。したがって、
\[
\frac{d}{dt}\hat{A}_H(t) = \frac{1}{i\hbar}[\hat{A}_H(t), \hat{H}] + \left(\frac{\partial \hat{A}}{\partial t}\right)_H
\]
を得る。これは古典力学におけるポアソン括弧を用いたハミルトンの運動方程式と直接的な対応関係にある。

\subsection{局所演算子としての場}
ハイゼンベルク描像では、時間に依存するのは状態ではなく演算子である。多自由度系や連続自由度系では、この考えを各空間点にラベル付けされた演算子族へ拡張できる。相対論的理論では、時間と空間をまとめて時空座標 $x^\mu=(t,\mathbf{x})$ とみなし、局所演算子 $\hat{\phi}(x)$ や $\hat{\psi}(x)$ を基本変数とする方が自然である。

この視点に立つと、粒子は一次的な存在ではなく、場の励起として理解される。次節では、この考えを最も単純な自由スカラー場で具体化する。

\section{場の量子論の基本枠組み}

\subsection{粒子ではなく場を量子化する}
前節では、相対論的理論では時空上の局所演算子を基本変数とするのが自然であることを見た。ここでは、その考えを具体的な量子理論として定式化する。

前節までの議論から分かるように、相対論的な一粒子波動方程式だけでは、負エネルギー解や反粒子の存在を無理なく扱うことができない。さらに、相対論と量子論を同時に考えると、十分高いエネルギーでは粒子対の生成・消滅が起こりうるため、粒子数を固定したヒルベルト空間そのものが根本的ではなくなる。

そこで場の量子論では、粒子を基本変数とみなすのではなく、時空の各点に値を割り当てる「場」を基本変数とする。記号 $x$ は時刻と位置をまとめた $x=(t,\mathbf{x})$ を表し、$\phi(x)$ や $\psi(x)$ は「時刻 $t$、位置 $\mathbf{x}$ における場の値」である。粒子とは、それらの場を量子化したときに現れる励起状態として理解される。こうすると、粒子の生成・消滅や反粒子の存在が、理論の枠組みの内部で自然に記述できる。

以下では、まず最も簡単な例である「自由な実スカラー場」を用いて、場の量子論の骨格を順に見る。「自由」とは他の場との相互作用がないこと、「スカラー」とはローレンツ変換の下で成分が混ざらず1つの値として変換することを意味する。

\subsection{自由スカラー場の古典論}
最も単純な例として、実スカラー場 $\phi(x)$ を考える。以下、この節では見通しをよくするために自然単位系 $c=\hbar=1$ を用いる。この単位系では、特殊相対論の式や量子論の式に現れる $c,\hbar$ をいちいち書かずに済む。
これまではまず古典論でエネルギーの表式を求めてから、その表式に現れる物理量を演算子で置き換えてきたのであった。今回もそのやり方に倣って話を進めてみよう。

場の理論でも、粒子力学と同様に、まず古典的な運動方程式を定めてから量子化する。その出発点になるのがラグランジアン密度 $\mathcal{L}$ である。これは各時空点に割り当てられた量であり、粒子力学のラグランジアン $L$ の場版と思えばよい。これを時空全体で積分した
\[
S = \int d^4x\, \mathcal{L}
\]
を作用という。ここで $d^4x = dt\,d^3x$ は時刻と空間全体にわたる積分を表す。作用が停留するという条件から運動方程式が決まる。

自由スカラー場のラグランジアン密度を
\[
\mathcal{L} = \frac{1}{2}\partial_\mu \phi \, \partial^\mu \phi - \frac{1}{2}m^2\phi^2
\]
とおく。ここで $\partial_\mu$ は時刻と空間に関する微分をまとめた記号であり、$\partial_\mu \partial^\mu$ は時間微分と空間微分をローレンツ対称な形でまとめた演算子である。作用 $S$ に対してオイラー--ラグランジュ方程式を適用すると、
\[
\partial_\mu \partial^\mu \phi + m^2 \phi = 0
\]
が得られる。これはすでに見たクライン--ゴルドン方程式そのものである。したがって、クライン--ゴルドン方程式は「一粒子波動方程式」としてだけでなく、「古典場の運動方程式」としても読めることが分かる。

次に、後でハミルトン形式で量子化するため、場に対する「共役な運動量」を定義する。これは粒子力学における座標 $q$ に対する運動量 $p$ に対応する量であり、
\[
\pi(\mathbf{x},t) := \frac{\partial \mathcal{L}}{\partial \dot{\phi}(\mathbf{x},t)} = \dot{\phi}(\mathbf{x},t)
\]
で与えられる。ここで $\dot{\phi}$ は $\phi$ の時間微分である。

同様に、粒子力学のハミルトニアンに対応する量としてハミルトニアン密度
\[
\mathcal{H} = \pi \dot{\phi} - \mathcal{L}
= \frac{1}{2}\pi^2 + \frac{1}{2}(\nabla \phi)^2 + \frac{1}{2}m^2\phi^2
\]
を得る。全ハミルトニアンはこれを空間全体で積分した
\[
H = \int d^3x\, \mathcal{H}
\]
である。したがって、自由場は「各空間点に自由度をもつ系」、あるいはフーリエ変換すれば「連続無限個の調和振動子の集まり」とみなせる。

\subsection{正準量子化}
ここから古典場を量子化する。正準量子化とは、粒子力学で $q,p$ を演算子 $\hat{q},\hat{p}$ に置き換えて交換関係 $[\hat{q},\hat{p}] = i$ を課したのと同じ考え方を、場 $\phi,\pi$ に適用する手続きである。

すなわち、古典場 $\phi,\pi$ を演算子値の場 $\hat{\phi},\hat{\pi}$ に持ち上げ、同時刻での正準交換関係
\[
[\hat{\phi}(\mathbf{x},t), \hat{\pi}(\mathbf{y},t)] = i\delta^{(3)}(\mathbf{x}-\mathbf{y}),
\]
\[
[\hat{\phi}(\mathbf{x},t), \hat{\phi}(\mathbf{y},t)] = 0, \quad
[\hat{\pi}(\mathbf{x},t), \hat{\pi}(\mathbf{y},t)] = 0
\]
を課す。ここで $\delta^{(3)}(\mathbf{x}-\mathbf{y})$ は3次元デルタ関数であり、「異なる空間点の自由度は独立で、同じ点でだけ基本交換関係が成り立つ」ことを表している。

これは、粒子力学における $[\hat{q},\hat{p}] = i$ の連続自由度版である。

\subsection{モード展開と生成・消滅演算子}
自由スカラー場の理論を考える時は、空間の各点ごとに独立な自由度がばらばらにあると考えるよりも、場全体がひとつの力学系を構成し、その振動モードが量子化されると理解すると見通しがよい。
自由場の方程式は線形であるから、その解は平面波の重ね合わせで表せる。量子化した後も同じ考え方が使え、演算子場は各運動量モードの和として
\[
\hat{\phi}(x) =
\int \frac{d^3p}{(2\pi)^3}\frac{1}{\sqrt{2E_{\mathbf{p}}}}
\left(
\hat{a}_{\mathbf{p}} e^{-ip\cdot x}
+ \hat{a}_{\mathbf{p}}^\dagger e^{ip\cdot x}
\right),
\qquad
E_{\mathbf{p}} = \sqrt{\mathbf{p}^2 + m^2}.
\]
と展開できる。これをモード展開という。指数関数 $e^{\mp ip\cdot x}$ は各運動量モードの時空依存性を表し、$E_{\mathbf{p}} = \sqrt{\mathbf{p}^2 + m^2}$ は相対論的エネルギー・運動量関係である。

ここで $\hat{a}_{\mathbf{p}}$ と $\hat{a}_{\mathbf{p}}^\dagger$ は、それぞれ運動量 $\mathbf{p}$ をもつ量子を一つ消滅・生成する演算子であり、
\[
[\hat{a}_{\mathbf{p}}, \hat{a}_{\mathbf{q}}^\dagger] = (2\pi)^3 \delta^{(3)}(\mathbf{p}-\mathbf{q}),
\quad
[\hat{a}_{\mathbf{p}}, \hat{a}_{\mathbf{q}}] = 0
\]
を満たす。この交換関係は、各運動量モードが調和振動子と同じ代数構造を持つことを示している。

ハミルトニアンはこれらを用いて書き直せる。すると各モードは「エネルギー $E_{\mathbf{p}}$ をもつ量子が何個あるか」で記述される形になる。真空にも形式的には無限大の零点エネルギーが現れるが、まずはその基準値を引いた量を見るため、生成演算子を左、消滅演算子を右に並べる正規順序をとると、
\[
:\hat{H}:=
\int \frac{d^3p}{(2\pi)^3} E_{\mathbf{p}} \,
\hat{a}_{\mathbf{p}}^\dagger \hat{a}_{\mathbf{p}}
\]
となる。したがって、一つ一つの粒子は「場の一つのモードを一量子だけ励起した状態」に対応している。

\subsection{フォック空間と多粒子状態}
粒子がひとつも存在しない状態を真空状態と呼び、$|0\rangle$ と書く。これは
\[
\hat{a}_{\mathbf{p}}|0\rangle = 0 \qquad (\forall \mathbf{p})
\]
で定義される。つまり、どの運動量モードについても「これ以上消せる粒子がない」状態である。

このとき、一粒子状態は $\hat{a}_{\mathbf{p}}^\dagger|0\rangle$、二粒子状態は $\hat{a}_{\mathbf{p}}^\dagger \hat{a}_{\mathbf{q}}^\dagger|0\rangle$ のように作られる。一般に
\[
|\mathbf{p}_1,\dots,\mathbf{p}_n\rangle
= \hat{a}_{\mathbf{p}_1}^\dagger \cdots \hat{a}_{\mathbf{p}_n}^\dagger |0\rangle
\]
が $n$ 粒子状態を与える。ここでは、以前に固定した粒子数 $n$ を最初から決めておく必要はなく、生成演算子を何回作用させたかによって粒子数が決まる。

このように、場の量子論で用いるヒルベルト空間は、粒子数が固定されたひとつの空間ではなく、$0$ 粒子、$1$ 粒子、$2$ 粒子、$\dots$ の各部分空間をすべて合わせた空間である。これをフォック空間という。相互作用を導入すると、時間発展の中でこれら異なる粒子数の部分空間が混ざり合う。ここに、粒子対生成や消滅を自然に扱える理由がある。

\subsection{ディラック場への一般化}
同じ考え方はスカラー場だけでなく、電子のようなスピン $1/2$ 粒子を記述するディラック場にも適用できる。ディラック方程式も、もはや一粒子の波動方程式としてではなく、スピノール場 $\psi(x)$ に対する古典場の方程式として読み直すのが本質的である。

ただし、電子はフェルミ粒子であるから、ボース粒子のような交換関係ではなく、同時刻の反交換関係
\[
\{\hat{\psi}_\alpha(\mathbf{x},t), \hat{\psi}_\beta^\dagger(\mathbf{y},t)\}
= \delta_{\alpha\beta}\delta^{(3)}(\mathbf{x}-\mathbf{y})
\]
を課す。ここで添字 $\alpha,\beta$ はスピノールの成分を表す。これにより、電子のようなフェルミ粒子に対して同じ量子状態を二つ以上占有できないというパウリの排他原理が自動的に組み込まれる。

ディラック場でもモード展開を行うと、電子を生成する演算子と陽電子を生成する演算子が別々に現れる。したがって、反粒子の存在はもはや後付けの再解釈ではなく、場の量子化そのものの帰結となる。これは、第6章で見た「負エネルギー解をどう解釈するか」という問題に対する、より根本的な答えである。

\subsection{相互作用への道筋}
自由場の量子化によって、粒子の概念が「場の励起」として再構成された。次に必要なのは、異なる場同士の相互作用を入れることである。

たとえば、電磁場 $A_\mu(x)$ とディラック場 $\psi(x)$ の最も基本的な相互作用は
\[
\mathcal{L}_{\mathrm{int}} = -e \bar{\psi}\gamma^\mu A_\mu \psi
\]
で記述される。ここで $e$ は電荷、$\bar{\psi} = \psi^\dagger \gamma^0$ はディラック共役と呼ばれる量であり、$\gamma^\mu$ はローレンツ対称性を保つための行列である。この式は、「電子場と電磁場が時空の各点で局所的に結合する」ということを表している。

これが量子電磁力学（QED）の出発点となる。相互作用を摂動的に扱うことで、散乱振幅、断面積、崩壊幅といった観測量がファインマン図を通じて計算される。

この意味で、場の量子論とは単に「相対論的な量子力学」ではなく、相対論・量子論・粒子数の変化・局所相互作用を同時に扱うための統一的な言語なのである。

\end{document}
